\scanpage{307}
\mspage{223}

Le foncteur (4.9.6.3) étant cofinal, pour établir que \unsure{4.9.2.1)} l'est, il suffit de prouver qu'il en est ainsi du foncteur (4.9.6.2), \inserted{ce que nous ferons.} Notons de plus que ledit foncteur est \uline{fidèle}, car il s'insère dans le carré commutatif

\diagram{
\begin{tikzcd}[ampersand replacement=\&]
\underline{\operatorname{Hom}}(J,C)_{/\varphi} \arrow[r] \arrow[d] \& \prod C_{/\varphi(j)} \arrow[d] \\
\underline{\operatorname{Hom}}(J,C) \arrow[r] \& \deleted{\unsure{\prod|F}} C^{\operatorname{Ob}J}
\end{tikzcd}
}

où les flèches verticales et la deuxième flèche horizontale sont fidèles. \deleted{\unsure{On va|En ce}} \deleted{alors} En définitive, pour établir le th. \unsure{4.9.3}, nous allons étudier le foncteur \inserted{\uline{fidèle}} (4.9.6.2)

\[
f\colon\underline{\operatorname{Hom}}(J,C)_{/\varphi}\unsure{\longrightarrow|\text{---}}\prod_{j\in J_0}C_{/\varphi(j)}\tag*{($\ast$)}
\]

(où on a posé $J_0=\operatorname{Ob}J$, $J_1=\mathrm{Fl}\,J$), et \deleted{il \unsure{faut|faudra} nous} nous allons prouver :

\textcircled{a} \deleted{\unsure{Le}}Ce foncteur \inserted{fidèle} $(\ast)$ est \uline{cofinal}.

\textcircled{b} La catégorie-source \inserted{$\underline{\operatorname{Hom}}(J,C)_{/\varphi}$} est gde devant $\pi_0$.

Nous allons pour ceci dégager un lemme \unsure{auxiliaire|critérium|condition} sur la cofinalité des foncteurs \uline{fidèles} :
% archive sheet number: 306
